\documentclass[11pt,a4paper]{article}
\usepackage[utf8]{inputenc}
\usepackage[T1]{fontenc}
\usepackage[portuguese]{babel}
\usepackage{lmodern}
\usepackage{microtype}
\usepackage{amsmath,amssymb}
\usepackage[a4paper,margin=2.3cm]{geometry}
\usepackage{xcolor}
\usepackage{titlesec}
\usepackage{enumitem}

\definecolor{azul}{HTML}{1E4E79}
\definecolor{verdee}{HTML}{16A34A}
\titleformat{\section}{\large\bfseries\color{azul}}{}{0em}{}
\newcommand{\slide}[3]{\medskip\noindent{\color{azul}\rule{3pt}{14pt}}\;
  \textbf{Slide #1 --- #2}\hfill\textit{(#3)}\\[1mm]}
\newcommand{\nucleo}{{\color{verdee}\footnotesize$\bigstar$\,núcleo}}
\setlength{\parskip}{4pt}
\setlength{\parindent}{0pt}

\title{\textbf{Roteiro da apresentação}\\[2mm]
\large Política monetária no Chile --- modelo Novo-Keynesiano mínimo,\\calibração, previsão e análise macro (Dynare)}
\author{Daniel Colli e João Zangrandi --- FGV EESP}
\date{\today}

\begin{document}
\maketitle

\noindent\textit{Como usar: o texto em corpo normal é o que se fala; os marcadores azuis indicam o
slide e o tempo. A apresentação completa tem 37 quadros e dura $\approx$28 minutos de fala. Para uma
versão de \textbf{20 minutos}, apresente os slides marcados com \nucleo{} (a \textbf{rota principal},
1--23 e 31) e trate a seção de \textbf{Análise Macro Aprofundada} (24--30) e o \textbf{Apêndice}
(32--37) como material para perguntas. Os números-chave estão em \textbf{negrito}.}

\bigskip

% ===================================================================
\section{Rota principal --- o modelo, a calibração e a previsão}
% ===================================================================

\slide{1--2}{Título e roteiro}{0,5 min · \nucleo}
Boa tarde. Nós vamos mostrar como construir e usar um \textbf{modelo de política monetária} para o
Chile. A ideia central é simples: um banco central tem basicamente \emph{um} instrumento, a taxa de
juros, e precisa decidir como mexê-la para controlar a inflação sem desestabilizar a economia. Vamos
começar pelo modelo e a intuição; depois calibrá-lo com dados reais do Chile; gerar previsões com ele;
e, no fim, aprofundar com testes e extensões.

\slide{3}{A pergunta da política monetária}{1 min · \nucleo}
Vamos pela pergunta. O Banco Central do Chile move a TPM, a taxa de política monetária --- o
equivalente à nossa Selic. Subir o juro encarece o crédito, esfria consumo e investimento, e tende a
derrubar a inflação; baixar faz o contrário. O objetivo é manter a inflação na \textbf{meta de 3\% ao
ano}, sem oscilações desnecessárias na atividade. Há um conflito embutido: às vezes, para baixar a
inflação, é preciso esfriar a economia. Para estudar isso, acompanhamos três variáveis: o
\textbf{hiato do produto} (se a economia está acima ou abaixo do potencial), a \textbf{inflação} e o
\textbf{juro}. Um detalhe: o modelo trabalha com elas em \emph{desvios} de um equilíbrio, então ficam
todas em torno de zero --- um truque que transforma equações complicadas em equações de linha reta.

\slide{4}{O modelo em três equações}{2 min · \nucleo}
Aqui está o coração de tudo: três equações. A primeira é a \textbf{curva IS}, a \emph{demanda}: ela diz
que a demanda hoje depende da demanda esperada amanhã e do \textbf{juro real} comparado ao \textbf{juro
neutro}. Juro real é o juro nominal menos a inflação esperada; o neutro, que chamamos de
\emph{r-estrela}, é o juro que deixa a economia equilibrada. Se o juro real está acima do neutro, a
política está apertada e a demanda cai. A segunda é a \textbf{curva de Phillips}, a \emph{oferta}: a
inflação de hoje depende da inflação esperada e do aquecimento da economia, com intensidade
\emph{kappa}. Guardem um ponto: essa curva só olha para frente, não tem inflação passada --- ou seja,
o modelo não tem ``inércia'' de inflação, o que vai importar lá na frente. A terceira é a \textbf{regra
de Taylor}, o comportamento do BC: ele sobe o juro quando a inflação sobe (intensidade \emph{phi-pi}) e
quando a economia aquece (\emph{phi-x}), e faz isso com \textbf{suavização} --- o parâmetro
\emph{rho-i} captura que bancos centrais mexem o juro em passos graduais. Resumindo: demanda reage ao
juro real; inflação reage ao aquecimento e ao futuro; e o BC reage à inflação e ao hiato, devagar.

\slide{5}{Calibração e a tabela parâmetro--alvo}{1 min · \nucleo}
``Calibrar'' é escolher os números do modelo, e é o que o trabalho pede para o Chile. Fixamos o juro
neutro em \textbf{3\% ao ano} --- testando também 2\% e 4\% --- e disso derivamos o fator de desconto
beta, \textbf{0,9926} (beta é um sobre um mais o juro neutro). A elasticidade sigma fica em 1; a
inclinação de Phillips kappa em \textbf{0,10} (testamos 0,07 e 0,13); a suavização rho-i nós
\emph{estimamos} dos dados e dá \textbf{0,934}, bem alto; phi-pi fica em 1,75 e phi-x em 0,5. Por fim,
os tamanhos dos choques são calibrados para que a decomposição de variância fique realista. Tudo isso
vira a tabela ``parâmetro contra alvo'', um dos produtos pedidos.

\slide{6}{Os dados do Chile}{1 min · \nucleo}
Estes são os dados, oficiais do Banco Central do Chile, trimestrais, de 2001 a 2026 --- 101
trimestres. Em cima, a TPM, que variou de menos de 1\% a quase 11\%. No meio, a inflação, com a linha
pontilhada na meta de 3\%: dá para ver os picos de 2008 e de 2022, acima de 10\%. Embaixo, o hiato,
com o tombo de \textbf{menos 15\%} na pandemia. O hiato sai do filtro de Hodrick--Prescott, que separa
tendência de ciclo. Um detalhe que importa: para alimentar o modelo, tiramos a média de cada série --
e a média da inflação no período foi \textbf{3,79\%}, acima da meta de 3\%; isso volta na previsão.

\slide{7}{Solução e determinação}{1 min · \nucleo}
Resolver o modelo é achar uma trajetória única e estável, apesar de as variáveis dependerem do futuro.
A ferramenta é a condição de \textbf{Blanchard--Kahn}. A regra: cada variável que ``salta'' --- aqui a
inflação e o hiato --- precisa de exatamente uma raiz explosiva para ser amarrada. Temos duas variáveis
de salto e um estado, o juro passado. O Dynare confirma duas raízes explosivas e uma estável, então a
solução é \textbf{única e estável}. O que garante isso é o \textbf{princípio de Taylor}: o BC tem de
reagir à inflação mais que proporcionalmente, phi-pi maior que 1; senão, surgem profecias
auto-realizáveis de inflação. Vamos ver isso desenhado.

\slide{8--9}{IRFs do baseline}{1,5 min · \nucleo}
Este é o gráfico mais importante: as \textbf{funções de resposta a impulso}. Elas respondem: se eu der
um choque hoje, o que acontece nos próximos trimestres? As linhas são hiato, inflação e juro; as
colunas, os três choques. Num \textbf{choque de demanda}, o hiato sobe, a inflação sobe e o BC reage
subindo o juro. Num \textbf{choque de custo} --- tipo petróleo --- a inflação sobe, mas o hiato
\emph{cai}: é o trade-off clássico, conter inflação custa atividade. Num \textbf{choque de política},
um aperto-surpresa, hiato e inflação caem, com resposta gradual pela suavização. Os sinais batem com a
teoria e tudo converge rápido.

\slide{10}{Momentos: modelo versus dados}{1 min · \nucleo}
Aqui comparamos as estatísticas do modelo com as dos dados. Duas lições. Primeira: o modelo
\emph{subestima} a volatilidade real --- faz sentido, a economia teve choques gigantes que um modelo
mínimo não reproduz. Segunda, e mais importante: a \textbf{persistência da inflação} nos dados é cerca
de 0,4, mas no modelo é só 0,05. É a confirmação numérica daquilo: como a curva de Phillips só olha
para frente, o modelo não tem inércia. É um limite conhecido e honesto.

\slide{11--12}{FEVD --- decomposição da variância}{1 min · \nucleo}
A decomposição da variância responde: de toda a oscilação de uma variável, quanto vem de cada choque? O
\textbf{hiato} vem de demanda e política, meio a meio; a \textbf{inflação} é dominada por choques de
\textbf{custo}, cerca de 80\%; o \textbf{juro}, pelo seu próprio choque, uns 75\%. É o padrão esperado
de um modelo Novo-Keynesiano, e o gráfico mostra que a calibração dos choques de fato atinge esses
alvos.

\slide{13--14}{Sensibilidade a kappa}{1 min · \nucleo}
Agora a sensibilidade, que o trabalho pede. Variando \textbf{kappa} entre 0,07 e 0,13: quanto maior,
mais forte o repasse do aquecimento para a inflação, e mais o BC precisa reagir. O segundo gráfico
resume o trade-off --- impacto, persistência e custo em produto --- para cada kappa.

\slide{15--16}{Sensibilidade a phi-pi}{1 min · \nucleo}
Variando \textbf{phi-pi} entre 1,3 e 2,2: um phi-pi maior reduz o impacto inicial de um choque de custo
sobre a inflação e acelera um pouco a convergência, mas exige movimentos maiores de juros e de produto.

\slide{17--18}{Mapa de determinação}{1 min · \nucleo}
Este é, na minha opinião, o gráfico mais bonito. Na faixa \textbf{vermelha}, phi-pi menor que 1: há
raízes estáveis demais --- indeterminação. Na faixa \textbf{verde}, phi-pi maior ou igual a 1: solução
única. Vejam a linha: a raiz dominante despenca de quase 1 para cerca de 0,66 no instante em que phi-pi
cruza 1, e continua caindo. Ou seja: quanto mais ativo o BC, mais rápida a convergência. É o princípio
de Taylor desenhado.

\slide{19}{Persistência: rho-i estimado vs calibrado}{0,5 min · \nucleo}
Comparando rho-i estimado, 0,934, com o calibrado, 0,80: com a suavização mais alta, o juro se move
mais devagar e os efeitos da política duram mais.

\slide{20}{Como a previsão é feita --- nativa no Dynare}{1 min · \nucleo}
Chegamos à previsão, a parte central, das páginas 48 a 55. E quero enfatizar: ela é feita
\emph{nativamente no Dynare}. A receita: declarar as variáveis observadas com \texttt{varobs} e rodar
\texttt{estimation} com \texttt{forecast = 8}. O detalhe crucial é \texttt{mode\_compute = 0}: isso diz
para \emph{não} estimar nada, manter a calibração fixa. O Dynare então roda o filtro de Kalman sobre o
histórico e, do último estado, projeta oito trimestres, devolvendo a previsão central e as bandas de
incerteza.

\slide{21}{Previsão incondicional --- fan chart}{1 min · \nucleo}
``Incondicional'' quer dizer: o modelo seguindo suas próprias leis, sem impor nada. A linha azul é a
projeção; a faixa, a incerteza. A inflação sai de 5,7\% e converge para perto de \textbf{3,8\%}, o juro
cai de 4,5\% para perto de \textbf{4,1\%} e o hiato fecha. Reparem no \emph{salto} da inflação: ela
estava em 5,7\% e a projeção começa em 3,3\%, porque o modelo não tem inércia e trata o pico como
transitório. É uma característica real, não um erro.

\slide{22}{Previsões condicionais --- cenários de política}{1 min · \nucleo}
``Condicional'' é o contrário: a gente \emph{impõe} uma trajetória a uma variável e o modelo ajusta o
resto, com o comando \texttt{conditional\_forecast}. Fizemos três cenários: manter a TPM em 4,5\%
(levemente contracionista, inflação cai a 3,0\%), apertar para 5,1\% (forte desinflação, mas o hiato
vai a menos 1,6\%) e forçar a inflação à meta de 3\% (exige aperto, o juro sobe).

\slide{23}{Choques implícitos}{0,5 min · \nucleo}
Para impor cada trajetória, o BC precisa empurrar o juro com choques de política. Este gráfico mostra o
tamanho desses empurrões --- torna explícito o esforço embutido em cada cenário.

% ===================================================================
\section{Aprofundamento --- análise macro (opcional na rota de 20 min)}
% ===================================================================

\slide{24}{Decomposição histórica}{1,5 min}
Aqui o modelo ``lê a história''. Normalmente ele roda para frente; agora roda \emph{para trás}, como um
detetive: dado o que aconteceu, o suavizador de Kalman infere quais choques ocorreram. As barras fatiam
a inflação observada em \textbf{demanda (azul)}, \textbf{custo (laranja)} e \textbf{política (verde)}. O
pico de 2022 é custo mais demanda; a COVID foi um colapso de demanda. Como há três observáveis e três
choques, o encaixe é exato por construção --- é a narrativa do modelo, não uma prova causal.

\slide{25}{SVAR versus DSGE}{1 min}
Aqui testamos se um método estatístico ``sem teoria'' concorda com o nosso modelo. Comparamos a
resposta a um choque de juros do DSGE e de um VAR estrutural. Eles discordam no curto prazo: o VAR
mostra o hiato \emph{subindo} após o aperto, o chamado \emph{activity puzzle}. Isso é honesto: mostra
que o resultado depende das suposições de identificação; o VAR desafia o modelo, não o valida.

\slide{26}{Economia aberta e inércia}{1 min}
Duas doses de realismo. À esquerda, abrimos a economia: câmbio, paridade de juros, \emph{pass-through}
de preços importados e um choque de cobre --- o Chile é exportador de commodity. À direita, demos
\emph{memória} à inflação com a curva de Phillips híbrida, que adiciona inflação passada: a
persistência sobe de 0,05 para \textbf{0,38}, perto dos dados. Ambos os modelos continuam determinados.

\slide{27}{MCMC --- parâmetros são incertos}{1 min}
Até aqui usamos um ``melhor palpite'' por parâmetro. O MCMC sorteia milhares de combinações compatíveis
com os dados --- duas cadeias de dez mil --- e mapeia a incerteza. À esquerda, a faixa de 90\% de cada
parâmetro; à direita, o R-chapéu de convergência, todos abaixo de 1,05. Achado honesto: a faixa de
phi-pi encosta abaixo de 1, ou seja, não dá para cravar que o BC reage forte o suficiente.

\slide{28}{Qual Phillips os dados preferem?}{1 min}
Entre a Phillips simples e a com inércia, qual os dados preferem? Usamos a evidência marginal, que
mede o ajuste já penalizando a complexidade. A \textbf{híbrida vence por 12 log-pontos}, um fator de
Bayes de cerca de \textbf{166 mil} --- os dados preferem fortemente a versão com inércia, confirmando
que a maior fraqueza do modelo básico era real.

\slide{29}{Resposta direta --- a previsão}{1,5 min}
Então, depois de tudo, qual a nossa aposta? Hoje a inflação está em 5,7\%, o juro em 4,5\% e a economia
levemente abaixo do potencial. A melhor previsão diz que a economia \textbf{normaliza em cerca de um
ano}: inflação para \textbf{3,7\%}, juro para \textbf{4,2\%} e hiato fechando. Eu destaco o horizonte
de \emph{um ano} de propósito --- foi onde o teste fora-da-amostra mostrou o modelo batendo previsões
ingênuas em 16 a 32\%. A linha tracejada vermelha é a previsão da híbrida: ela segura a inflação mais
alta no curto prazo, mas converge para a mesma previsão de um ano.

\slide{30}{A previsão é robusta? O teste da híbrida}{1,5 min}
Aqui fechamos a pergunta ``o baseline é mesmo o melhor?''. À esquerda, a inércia: a inflação da híbrida
decai bem mais devagar que a do baseline. À direita, o resultado do teste fora-da-amostra: a híbrida
\textbf{corta o erro da inflação de um trimestre em 19\%} --- a inércia cura a fraqueza de curto prazo
--- mas em um ano os dois \textbf{empatam}. Conclusão: a previsão de um ano é \textbf{robusta} à escolha
do modelo; para o próximo trimestre, a híbrida é a ferramenta melhor.

% ===================================================================
\section{Fechamento}
% ===================================================================

\slide{31}{Obrigado}{0,25 min · \nucleo}
Este slide é apenas o encerramento visual da apresentação. A síntese oral deve ser feita no fim do
slide 30: a previsão de um ano é robusta entre baseline e NKPC híbrida; o modelo básico é útil, mas a
inércia inflacionária melhora o curto prazo. Depois disso, encerre com ``obrigado'' e abra para
perguntas.

% ===================================================================
\section{Apêndice --- material para perguntas (slides 32--37)}
% ===================================================================

\slide{32}{Choques recuperados}{para perguntas}
São as três séries de choques que o suavizador reconstruiu, com picos enormes na COVID e em 2022. É o
``input'' da decomposição do slide 24.

\slide{33}{Contrafactual com phi-pi = 2,5}{para perguntas}
``E se o BC tivesse sido mais agressivo?''. As linhas quase se sobrepõem --- porque a suavização alta
faz o juro andar devagar de qualquer jeito, então mudar phi-pi muda pouco. O gradualismo domina.

\slide{34}{r-estrela variável no tempo}{para perguntas}
Em vez de fixar o juro neutro em 3\%, estimamos ele variando; termina perto de 0,85\%, sugerindo que o
neutro caiu. É uma proxy estatística, não uma estimação estrutural.

\slide{35}{Previsão fora-da-amostra}{para perguntas}
O teste de previsão de verdade, com a híbrida incluída. O baseline é pior em um trimestre, mas vence o
passeio aleatório em um ano; a híbrida melhora o curto prazo em 19\%. É o teste que valida a
resposta-direta.

\slide{36}{Dados externos: peso e cobre}{para perguntas}
As séries novas do modelo aberto: o câmbio peso/dólar e o preço do cobre, dados reais, mostrando o peso
enfraquecendo e o ciclo do cobre.

\slide{37}{Modelo aberto versus fechado}{para perguntas}
A resposta a um choque de juros é quase igual nos dois modelos: abrir a economia adiciona o canal de
câmbio, mas não muda dramaticamente a transmissão monetária básica neste modelo pequeno.

\end{document}
